% NOTATION
%   G                 -- generic undirected graph on vertex set V
%   G^*               -- true (data-generating) conditional independence graph
%   \mathcal{D}_p     -- space of decomposable graphs on p vertices
%   \mathcal{M}(G^*)  -- set of minimal triangulations of G^*
%   \text{PIP}_{ij}   -- Stage-1 posterior edge inclusion probability for edge (i,j)
%   \tau, r           -- HIW concentration and prior edge inclusion probability

Applied researchers routinely ask which pairs of variables remain associated after conditioning on the rest of their data. A macroeconomist studying financial spillovers needs to separate the country pairs that transmit shocks directly from those that move together only through third parties. In genomics, the goal is the same kind of separation between genes that regulate each other and genes that merely covary through a common upstream signal. A central banker reading a cross-country correlation matrix faces the related question of which associations survive after partialling out the rest of the panel. Gaussian graphical models (GGMs) are a natural tool for such questions.

The menu of GGM estimators is large and expanding. The graphical lasso \parencite{friedman2008sparse} and its semiparametric extension, the nonparanormal \parencite{liu2009nonparanormal}, return sparse point estimates selected by cross-validation. Bayesian approaches place a prior on the graph and return a full posterior, with a closed-form marginal likelihood when the graph is decomposable \parencite{dawid1993hyper,lauritzen1996graphical} and costly approximations when it is not \parencite{mohammadi2015bdgraph}. Each method carries its own sparsity assumption, its own tuning protocol, and its own implicit definition of the ``right'' graph. An applied researcher picking among them has little guidance on which method to use in which setting, how to interpret the output when the assumptions fail, or what the resulting object represents in economic terms.

This project engages with a specific strand of Bayesian GGM estimation. We focus on the empirical Bayes SAEM-MCMC procedure of \textcite{donnet2012empirical}. The procedure has two features applied researchers should value: data-driven regularization in place of a cross-validated tuning parameter, and a full posterior over graphs that supports formal uncertainty statements. Its main limitation is also clear: the search is restricted to decomposable graphs, a class that may exclude structures encountered in applied work.

Two questions follow. First, what does the decomposable posterior concentrate on when the true graph is non-decomposable, and is the resulting object useful? Second, can the structured bias of the decomposable posterior be converted into a data-driven prior for inference over the full graph space? \textcite{niu2020bayesian} answer the first question: under mild regularity, the posterior mode concentrates on the set of minimal triangulations $\mathcal{M}(G^*)$ of the true graph, missing a true edge is penalized exponentially in $n$, and including a spurious fill edge is penalized only polynomially. The decomposable posterior is therefore a structured over-estimate of $G^*$: it recovers true edges with high confidence but contaminates them with fill edges whose posterior inclusion probabilities drift toward $1/2$ when $G^*$ admits competing triangulations.

Our contribution is a two-stage procedure that exploits this characterization. Stage 1 runs SAEM-MCMC on the space of decomposable graphs $\mathcal{D}_p$ and returns the posterior inclusion probability (PIP) matrix, whose entries $\text{PIP}_{ij}$ give the posterior probability that edge $(i,j)$ is present. Stage 2 feeds the PIP matrix into the BDMCMC sampler over the unrestricted graph space \parencite{mohammadi2015bdgraph} through what we call a \textit{Decomposable-Informed Prior} (DIP),
\[
    p_{ij} \;=\; \lambda\,\text{PIP}_{ij} \;+\; (1-\lambda)\, r_0,
\]
where $r_0$ is a flat reference probability and $\lambda \in [0,1]$ is a mixing weight. Setting $(\lambda, r_0) = (0, 0.2)$ recovers the BDMCMC default; setting $\lambda > 0$ folds the Stage-1 PIPs into the Stage-2 prior. The procedure delivers two graphs from the same dataset: a decomposable Stage-1 graph $\hat G_D$ and an unrestricted Stage-2 graph $\hat G$. We evaluate the procedure on simulated Gaussian data covering five topologies (chain, star, cycle, sparse Erd\H{o}s--R\'enyi, and Barab\'asi--Albert) and on the daily volatility series of \textcite{demirer2018estimating} for $96$ international banks and $10$ sovereign bonds over 2003--2014. The decomposability assumption is consequential but not always costly. On topologies whose minimal triangulation is unique or near-unique, the deviation from unrestricted Bayesian estimation is small. Combined with the closed-form marginal likelihood and the absence of normalizing-constant approximations, this makes the decomposable family a sensible default at moderate $p$, even when the truth is technically not chordal. 
%at $p \in \{5, 10, 20\}$ and four signal strengths $\rho_{\max} \in \{0.13, 0.19, 0.38, 0.63\}$, and on the daily volatility series of \textcite{demirer2018estimating} for $96$ international banks and $10$ sovereign bonds over 2003--2014. Unsurprisingly, the BDMCMC sampler over the unrestricted graph space recovers more structure than the decomposable Stage-1 alone, by a mean $\Delta\text{MCC} \approx 0.13$ across the design. More usefully, the calibrated two-stage estimator tracks the BDMCMC default within $0.06$ MCC at every cell of the grid, which suggests that the decomposability restriction is not an extreme approximation on the designs we consider.
We make two contributions. First, we synthesize the \textcite{niu2020bayesian} characterization of posterior concentration on minimal triangulations into an applied-friendly account, with simulations that probe its finite-sample behavior. Second, we propose the DIP as a diagnostic, not a superior estimator, and show how the two-graph output reveals whether the data support a decomposable or non-decomposable structure. 

% We make two contributions. First, we synthesize the \textcite{niu2020bayesian} characterization into an applied-reader-friendly account of what the decomposable empirical Bayes posterior recovers when the truth is non-decomposable, with simulations that probe its applicability. Second, we propose the DIP and the associated two-stage estimator, whose two-graph output is a built-in misspecification diagnostic: the comparison between the Stage-1 decomposable graph and the Stage-2 unrestricted graph tells the practitioner whether the recovered structure lies in the regime where Theorem~\ref{thm:concentration} grants asymptotic concentration on $\mathcal{M}(G^*)$. We illustrate the diagnostic on the DDLY bank-volatility data in Section~\ref{sec:application}, where the recovered Stage-2 graph is markedly non-decomposable and the diagnostic restricts the qualitative claims to the parts that survive a decomposable cover.

Section~\ref{sec:theory} reviews decomposable graphs and the Hyper-Inverse Wishart marginal likelihood. Section~\ref{sec:algorithm} describes the SAEM-MCMC algorithm and our implementation improvements. Section~\ref{sec:misspec} develops the misspecification theory and introduces the DIP. Section~\ref{sec:sim} reports simulation evidence on decomposable and non-decomposable truths. Section~\ref{sec:application} presents the empirical application. Section~\ref{sec:discussion} discusses the results and suggests next steps toward a practitioner guide.
